% !PRINTED 377
\weberpage{377}
\begin{center}
§. 122.

Quadratische Irrationalzahlen.
\end{center}

Einfache und schöne Gesetze ergeben sich, wenn man die Kettenbruchentwickelung auf die Bestimmung der Wurzeln einer quadratischen Gleichung anwendet. Die Wurzel einer ganzzahligen quadratischen Gleichung hat die Form
\weq{1}{\omega = x + y\sqrt{d},}
worin $x$, $y$, $d$ ganze oder gebrochene rationale Zahlen sind. Wir können aber, ohne die Allgemeinheit zu beschränken, $d$ als \gesperrt{ganze Zahl und ohne quadratischen Theiler voraussetzen}; denn wenn $d$ einen Nenner hat, so können wir mit diesem Nenner erweitern und können die Wurzel aus dem quadratischen Nenner und aus einem etwaigen quadratischen Theiler des Zählers zu $y$ rechnen; wir nehmen dann weiter an, dass $d$ nicht gleich 1 ist, da sonst $\omega$ rational wäre; $d$ kann positiv oder negativ sein, und davon hängt es ab, ob $\omega$ reell oder imaginär ist.

Aendern wir das Vorzeichen der Wurzel, so erhalten wir
\weq{2}{\omega' = x - y\sqrt{d},}
was die zu $\omega$ \gesperrt{conjugirte} Zahl genannt wird. Das Product der beiden Zahlen $\omega$, $\omega'$
\[
\omega\omega' = x^2 - y^2 d
\]
heisst die \gesperrt{Norm von} $\omega$ (oder auch von $\omega'$). \gesperrt{Die Norm eines Productes zweier quadratischer Irrationalzahlen ist gleich dem Product der Normen.}

Die Zahlen $\omega$ und $\omega'$ sind die Wurzeln einer Gleichung mit rationalen Co\"efficienten:
\weq{3}{\omega^2 - 2x\omega + (x^2 - y^2 d) = 0.}
Um sie in eine ganzzahlige Gleichung zu verwandeln, setzen wir
\weq{4}{x^2 - y^2 d = -\frac{a}{c}, \quad 2x = \frac{b}{c},}
worin $a$, $b$, $c$ ganze Zahlen ohne gemeinsamen Theiler sind, und erhalten aus (3)
\weq{5}{c\omega^2 = a + b\omega,}
also eine ganzzahlige \gesperrt{primitive} Gleichung\par % [sic] source omits terminal punctuation; the next leaf begins a new sentence
